\section{Coin models}

Coin tossing is the simplest family of probability models, but it is not
trivial. It already contains several ideas that will appear throughout the
course: a sample space, a recording rule, a probability assignment, repeated
experiments, product spaces, events described by conditions, and the difference
between a particular sequence and a numerical summary of that sequence.

The physical object called a coin does not determine a probability model by
itself. We must say what is recorded and what assumptions are made about the
mechanism. A coin may be treated as symmetric, or it may be treated as biased.
A single toss may be recorded, or a whole sequence of tosses may be recorded.
Each of these choices changes the model or the probabilities inside the model.

\subsection*{One toss of a coin}

Suppose first that a coin is tossed once and that we record which side lands up.
The elementary outcome is one of two symbols:
\[
H \quad\text{or}\quad T,
\]
where \(H\) denotes heads and \(T\) denotes tails. The sample space is
\[
\Omega=\{H,T\}.
\]

At this point we have only described the possible outcomes. We have not yet
assigned probabilities.

If the coin is modelled as symmetric, then the two elementary outcomes are
given equal probability:
\[
P(H)=P(T)=\frac12.
\]
This is a modelling assumption based on symmetry. It is not a consequence of the
fact that the set \(\Omega\) has two elements. The fact that \(|\Omega|=2\) tells
us how many elementary outcomes there are. It does not by itself tell us how
likely they are.

If the coin is biased, we may instead write
\[
P(H)=p,\qquad P(T)=1-p,
\]
where
\[
0\leq p\leq 1.
\]
The same sample space \(\{H,T\}\) now carries a different probability
assignment. When \(p=1/2\), the model is symmetric. When \(p\neq 1/2\), one side
is more likely than the other.

\begin{remarkbox}
The symbols \(H\) and \(T\) are not probabilities. They are elementary outcomes.
The numbers assigned to them come from an additional assumption about the coin
and the tossing mechanism.
\end{remarkbox}

\subsection*{Events in the one-toss model}

In the one-toss model there are only four events:
\[
\varnothing,\qquad \{H\},\qquad \{T\},\qquad \Omega.
\]
Their probabilities are determined by the probabilities of \(H\) and \(T\):
\[
P(\varnothing)=0,\qquad P(\Omega)=1,
\]
\[
P(\{H\})=p,\qquad P(\{T\})=1-p.
\]
For a symmetric coin this becomes
\[
P(\{H\})=P(\{T\})=\frac12.
\]

Even this small example illustrates the general finite rule. If an event is a
set of elementary outcomes, then its probability is obtained by adding the
probabilities of the elementary outcomes inside it.

\subsection*{Two tosses: sequences as elementary outcomes}

Now suppose the coin is tossed twice and the complete order of results is
recorded. An elementary outcome is no longer a single symbol. It is an ordered
sequence of length two. The sample space is
\[
\Omega=\{HH,HT,TH,TT\}.
\]
The outcomes \(HT\) and \(TH\) are different because the first position and the
second position are different. The outcome \(HT\) means heads first and tails
second. The outcome \(TH\) means tails first and heads second.

If the coin is symmetric and the two tosses are modelled as independent
repetitions of the same mechanism, then every sequence has probability
\[
\frac14.
\]
Thus
\[
P(HH)=P(HT)=P(TH)=P(TT)=\frac14.
\]

If the coin has \(P(H)=p\) and \(P(T)=1-p\), then the probabilities of the four
sequences are
\[
P(HH)=p^2,\qquad P(HT)=p(1-p),
\]
\[
P(TH)=(1-p)p,\qquad P(TT)=(1-p)^2.
\]
The middle two values are equal because multiplication of numbers is
commutative:
\[
p(1-p)=(1-p)p.
\]
But the sequences \(HT\) and \(TH\) remain different elementary outcomes.

\begin{examplebox}
Let \(p=0.7\). Then
\[
P(HH)=0.49,\qquad P(HT)=0.21,\qquad P(TH)=0.21,\qquad P(TT)=0.09.
\]
These probabilities add to one:
\[
0.49+0.21+0.21+0.09=1.
\]
The sample space has four outcomes, but the model is not uniform.
\end{examplebox}

\subsection*{The event ``exactly one head''}

The event
\[
A=\{\text{exactly one head occurs}\}
\]
contains two elementary outcomes:
\[
A=\{HT,TH\}.
\]
Therefore, for a symmetric coin,
\[
P(A)=P(HT)+P(TH)=\frac14+\frac14=\frac12.
\]
For a biased coin with \(P(H)=p\),
\[
P(A)=p(1-p)+(1-p)p=2p(1-p).
\]

Notice the difference between an elementary outcome and an event. The sequence
\(HT\) is one elementary outcome. The phrase ``exactly one head'' describes an
event containing two elementary outcomes. Probability problems often ask for the
probability of an event, not of a single elementary outcome.

\subsection*{Three tosses and the structure of a product space}

For three tosses, with order recorded, the sample space is
\[
\Omega=\{H,T\}^3.
\]
This notation means the set of all ordered triples whose entries are \(H\) or
\(T\). Explicitly,
\[
\Omega=\{HHH,HHT,HTH,HTT,THH,THT,TTH,TTT\}.
\]
There are
\[
2^3=8
\]
elementary outcomes, because each of the three positions has two possible
symbols.

For a symmetric coin, each sequence has probability
\[
\frac18.
\]
For a biased coin with \(P(H)=p\), the probability of a sequence is obtained by
multiplying the probabilities of its symbols. For example,
\[
P(HTH)=p(1-p)p=p^2(1-p),
\]
and
\[
P(TTH)=(1-p)(1-p)p=(1-p)^2p.
\]

The exponent records how many times a symbol appears in the sequence. A sequence
with \(k\) heads and \(3-k\) tails has probability
\[
p^k(1-p)^{3-k}.
\]

\subsection*{The general model of \(n\) tosses}

For \(n\) tosses, with order recorded, the sample space is
\[
\Omega=\{H,T\}^n.
\]
An elementary outcome is a sequence
\[
\omega=(a_1,a_2,\ldots,a_n),
\]
where each \(a_i\) is either \(H\) or \(T\). There are
\[
|\Omega|=2^n
\]
such sequences.

If the coin is symmetric and each toss is modelled as an independent repetition,
then the model is uniform:
\[
P(\omega)=\frac1{2^n}
\]
for every \(\omega\in\Omega\).

If \(P(H)=p\) and \(P(T)=1-p\), then a sequence with \(k\) heads and \(n-k\)
tails has probability
\[
p^k(1-p)^{n-k}.
\]
This formula is not yet the probability of ``exactly \(k\) heads''. It is the
probability of one particular sequence with \(k\) heads.

\begin{examplebox}
For five tosses, the sequence
\[
HHTTH
\]
has three heads and two tails. Its probability in the biased model is
\[
p^3(1-p)^2.
\]
The sequence
\[
HTHTH
\]
also has three heads and two tails, so it has the same probability
\[
p^3(1-p)^2.
\]
But the two sequences are different elementary outcomes.
\end{examplebox}

\subsection*{Exactly \(k\) heads in \(n\) tosses}

Now consider the event
\[
A_k=\{\omega\in\{H,T\}^n: \omega\text{ contains exactly }k\text{ heads}\}.
\]
This event contains all sequences of length \(n\) in which exactly \(k\) positions
are occupied by \(H\), and the remaining \(n-k\) positions are occupied by \(T\).

To count these sequences, we do not choose the heads themselves as physical
objects. We choose the positions at which heads occur. There are \(n\) positions,
and we choose \(k\) of them to contain \(H\). Once those positions are chosen,
the remaining positions must contain \(T\). Hence
\[
|A_k|=\binom{n}{k}.
\]

For a symmetric coin, the whole space is uniform, so
\[
P(A_k)=\frac{|A_k|}{|\Omega|}=\frac{\binom{n}{k}}{2^n}.
\]

For a biased coin, the space is generally not uniform. However, every sequence
in \(A_k\) has the same probability \(p^k(1-p)^{n-k}\). Since there are
\(\binom{n}{k}\) such sequences, additivity gives
\[
P(A_k)=\binom{n}{k}p^k(1-p)^{n-k}.
\]

\begin{theorembox}
If a coin with \(P(H)=p\) is tossed \(n\) times independently, then the
probability of exactly \(k\) heads is
\[
\binom{n}{k}p^k(1-p)^{n-k},
\qquad k=0,1,\ldots,n.
\]
\end{theorembox}

This formula has three parts:
\begin{itemize}
    \item \(\binom{n}{k}\) counts the positions of the heads;
    \item \(p^k\) is the contribution of the \(k\) heads;
    \item \((1-p)^{n-k}\) is the contribution of the \(n-k\) tails.
\end{itemize}

\subsection*{What this model teaches}

The coin model teaches several lessons that will recur in more complicated
examples.

First, the sample space depends on what is recorded. One toss gives
\(\{H,T\}\); \(n\) ordered tosses give \(\{H,T\}^n\); recording only the number of
heads gives \(\{0,1,\ldots,n\}\).

Second, the probability assignment is separate from the sample space. The same
space \(\{H,T\}\) can carry the symmetric assignment \(1/2,1/2\) or the biased
assignment \(p,1-p\).

Third, an event may contain many elementary outcomes. The event ``exactly
\(k\) heads'' is not one sequence; it is a set of sequences.

Fourth, combinatorics enters probability when we count how many elementary
outcomes satisfy a condition. The binomial coefficient appears because we count
positions of successes, not because it is a formula to be memorized in isolation.

\begin{remarkbox}
The coin model is simple enough to be completely transparent, but rich enough to
show the entire modelling workflow: choose the recording rule, construct
\(\Omega\), assign probabilities, define events, count favourable outcomes, and
interpret the result.
\end{remarkbox}
